% -------------------------------------------------
% VERSION TRACKER — No modificar este bloque
% Versión:    Tesis Humanizada (Dummies)
% Archivo:    Capitulo0/Resumen.tex
% Última mod:  2026-05-06T16:15:31
% Git commit:  cc9c766f @ main
% Audiencia:   divulgativo
% Verificar antes de editar: bash check_tesis_version.sh overleaf_tesis/Capitulo0/Resumen.tex
% -------------------------------------------------


El presente trabajo de grado aborda el diseño y desarrollo de Pymetory, un sistema web para la gestión de inventario de materia prima orientado a pequeñas y medianas empresas (PYMEs) \cite{bbva2024mipymes,conpes4069} que actualmente enfrentan deficiencias en el control de insumos debido a la carencia de herramientas digitales especializadas. El proyecto parte de una entrevista de levantamiento de requisitos con la administradora de una panificadora del Valle del Cauca, cuyas necesidades reales ---diferencias entre el inventario contable y el físico, conciliación manual, reabastecimiento sin datos de consumo y trazabilidad por lote exigida por la normativa sanitaria--- guiaron los requerimientos del sistema junto con las cuatro reglas de negocio establecidas por la dirección del trabajo de grado.

El sistema fue implementado utilizando un stack tecnológico moderno y de código abierto: Laravel 13 (PHP 8.3) como framework backend bajo el patrón Modelo-Vista-Controlador (MVC), React 19 con TypeScript como capa de presentación, Inertia.js 2.0 como puente de comunicación entre backend y frontend, Tailwind CSS v4 para los estilos, y MySQL 8.0 como sistema gestor de base de datos relacional. La innovación principal del proyecto radica en la integración de un módulo de consulta inteligente basado en Large Language Models (LLM) con arquitectura RAG (Retrieval-Augmented Generation), que permite a los usuarios interactuar con el sistema mediante lenguaje natural en español colombiano y recibir respuestas construidas a partir de los datos reales del inventario, lo que reduce el riesgo de que el modelo de lenguaje invente cifras. El módulo RAG clasifica las consultas en 14 categorías de intención y recupera contexto pertinente desde las tablas de la base de datos antes de enviar el prompt al LLM; en una batería de 50 consultas verificadas contra la base de datos respondió correctamente las 50 (34 antes de las correcciones), y 8 de 10 en una prueba ciega con preguntas nuevas.

El sistema implementa el estándar FEFO (First Expired, First Out) como regla obligatoria de despacho, un Kardex digital inmutable con trazabilidad completa de cada movimiento, control de acceso basado en roles (RBAC) con perfiles de Admin y Operario, módulo de conciliación de inventario físico contra lógico, generación de reportes en PDF, CSV y Excel, etiquetas QR para identificación de lotes, un dashboard con indicadores clave en tiempo real y una proyección de reabastecimiento con punto de reorden calculada desde el Kardex. El proyecto se gestionó con un enfoque ágil inspirado en SCRUM, adaptado a un único desarrollador, sobre un tablero de GitHub Projects. El sistema cuenta con 178 pruebas PHPUnit (177 aprobadas y una omitida porque depende de un proveedor externo) y 61 pruebas de extremo a extremo con Playwright, todas aprobadas; la evaluación heurística de Nielsen corrigió 13 de 14 hallazgos, la auditoría WCAG 2.1 AA no encontró incumplimientos automáticos, y la prueba con usuario queda pendiente. El sistema se desplegó en la capa gratuita de Oracle Cloud Infrastructure y está disponible en \url{https://app.pymetory.com}.

El proyecto se enmarca en la modalidad de Trabajo de Investigación \cite{univalle_resolucion137}. Como resultado, se obtuvo una solución de software funcional, desplegada, verificada con pruebas automatizadas y de código abierto, que además de resolver un problema práctico, genera un modelo de referencia documentado sobre la integración de inteligencia artificial conversacional en sistemas de gestión empresarial para PYMEs.

\textbf{Palabras clave:} Gestión de inventarios, sistema web, Laravel, PYMEs, FEFO, Kardex digital, RAG (Retrieval-Augmented Generation), Large Language Models, investigación aplicada, logística de bodegas.
